% !Mode:: "TeX:UTF-8"

\hitsetup{
  %******************************
  % 注意：
  %   1. 配置里面不要出现空行
  %   2. 不需要的配置信息可以删除
  %******************************
  %
  %=====
  % 秘级
  %=====
  statesecrets={公开},
  natclassifiedindex={TM301.2},
  intclassifiedindex={62-5},
  %
  %=========
  % 中文信息
  %=========
  ctitleone={高动态连续旋转平面绳驱并联机器人},%本科生封面使用
  ctitletwo={关键技术研究},%本科生封面使用
  ctitlecover={高动态连续旋转平面绳驱并联机器人\\关键技术研究},%放在封面中使用，自由断行
  ctitle={高动态连续旋转平面绳驱并联机器人关键技术研究},%放在原创性声明中使用
  % csubtitle={一条副标题}, %一般情况没有，可以注释掉
  % cxueke={工学},
  csubject={机器人工程},
  caffil={机器人与先进制造学院},
  cauthor={柳尊文},
  csupervisor={袁晗教授},
  % cassosupervisor={某某某教授}, % 副指导老师
  % ccosupervisor={某某某教授}, % 联合指导老师
  % 如果是第一封面的日期要手动设置，需要取消注释下一行，并将内容改为“规范”中要求的封面第一页最下方的日期
  firstpagecdate={2026年5月},
  % 日期自动使用当前时间，若需指定按如下方式修改：
  cdate={2026年5月},
  cstudentid={220310426},
  % cstudenttype={学术学位论文}, %非全日制教育申请学位者
  % cnumber={no9527}, %编号
  % cpositionname={哈铁西站}, %博士后站名称
  % cfinishdate={20XX年X月---20XX年X月}, %到站日期
  % csubmitdate={20XX年X月}, %出站日期
  % cstartdate={3050年9月10日}, %到站日期
  % cenddate={3090年10月10日}, %出站日期
  %（同等学力人员）、（工程硕士）、（工商管理硕士）、
  %（高级管理人员工商管理硕士）、（公共管理硕士）、（中职教师）、（高校教师）等
  %
  %
  %=========
  % 英文信息
  %=========
  etitle={Research on Key Technologies for High-Dynamic Continuous Rotation Planar Cable-Driven Parallel Robots},
  esubtitle={This is the sub title},
  exueke={Engineering},
  esubject={Computer Science and Technology},
  eaffil={\emultiline[t]{School of Mechatronics Engineering \\ Mechatronics Engineering}},
  eauthor={Yu Dongmei},
  esupervisor={Professor XXX},
  eassosupervisor={XXX},
  % 日期自动生成，若需指定按如下方式修改：
  edate={December, 2017},
  estudenttype={Master of Art},
  %
  % 关键词用“英文逗号”分割
  ckeywords={绳驱并联机器人, 连续旋转, 力控制, 张力分配, 加速度跟踪},
  ekeywords={cable-driven parallel robot, continuous rotation, force control, tension distribution, acceleration tracking},
}

\begin{cabstract}

面向载人航天训练与复杂运动环境模拟需求，地面加速度模拟装备需要在较大工作空间内实现持续、高带宽且平滑的动态输出，以支持航天员在交会对接、姿态调整和长时间机动任务中的前庭适应与操控能力训练。传统离心机虽然能够提供较大的径向过载，但设备体积大、运行成本高，且运动形式相对单一，难以模拟多姿态耦合运动下的复杂加速度环境。绳驱并联机器人采用柔性绳索替代刚性连杆，具有工作空间大、动平台惯量低、负载自重比高和动态响应快等优点，为构建新型高动态加速度模拟平台提供了可行方案。

本文采用平面四绳冗余驱动方案构建机器人总体机构，完成了末端平台、出绳方式、关键锚点布局以及传感器配置的设计分析。针对传统平面出绳滑轮方案在包络长度补偿与结构复杂度方面的不足，本文提出了动态出绳点设计思路，以简化机构设计。该方案在后续样机迭代中完成了系统级设计，但本文算法验证主要基于力控制实验样机开展，其有效性不依赖于动态出绳点机构的实际应用。末端平台采用创新绕绳方法，实现了连续多圈旋转的效果。并结合张力传感器布置方案，为后续张力闭环控制与实验验证提供硬件基础。围绕机构控制需求，建立了绳长、平台位姿与驱动端状态之间的运动学关系，采用基于雅可比矩阵迭代的正运动学求解方法实现末端位姿估计，并在样机平台上完成位置控制流程验证。研究与仿真表明，基于逆运动学离散求解的常规位置控制能够实现给定轨迹的基本跟随，但在轨迹参数离散、绳索柔性、摩擦以及系统耦合等因素共同作用下，末端加速度容易出现不连续和抖振现象，难以满足高动态平滑加速度模拟要求。因此，本文将控制重点由位置控制转向力控制，尝试通过张力直接调节末端广义力，实现更连续的加速度输出。

在控制算法方面，本文围绕张力分配问题开展了分析、仿真与优化。针对平面四绳冗余驱动机构中末端期望加速度对应多组可行绳力解的特点，分析了正张力约束、力分布均衡性以及峰值载荷控制等要求对张力分配结果的影响，并结合轨迹生成过程对末端加速度和理论绳索力进行了仿真计算，为后续控制器设计提供了依据。

综合理论分析、仿真研究与阶段性实验结果可以看出，所设计的高动态连续旋转平面绳驱并联机器人方案具备开展连续旋转与高动态控制研究的基础，采用力控制代替单纯位置控制是改善末端加速度连续性、提升动态响应品质的合理技术路线。本文完成的平面四绳机构设计、运动学求解、位置控制验证、张力分配分析、力控制实验平台搭建以及单电机闭环力矩控制、四电机协同张力控制、末端平台加速度跟踪实验。本文研究可为连续旋转绳驱并联机器人的设计与控制提供参考，并为面向航天训练、飞行模拟及其他高动态运动模拟场景的装备研制提供理论依据与工程借鉴。
\end{cabstract}

\begin{eabstract}

Facing the requirements of manned spaceflight training and complex motion-environment simulation, ground-based acceleration simulation equipment is expected to provide continuous, high-bandwidth, and smooth dynamic output within a relatively large workspace so as to support astronauts' vestibular adaptation and manipulation training during rendezvous and docking, attitude adjustment, and long-duration maneuvering tasks. Traditional centrifuges can generate considerable radial overload, but they suffer from large size, high operating cost, and relatively limited motion patterns, which makes them unsuitable for reproducing complex acceleration environments under multi-attitude coupled motions. Cable-driven parallel robots, which replace rigid links with flexible cables, offer advantages such as large workspace, low moving-platform inertia, high payload-to-weight ratio, and fast dynamic response, and therefore provide a feasible solution for a new type of high-dynamic acceleration simulation platform.

In this work, a planar four-cable redundantly actuated configuration is adopted for the robot mechanism. The design and analysis of the moving platform, cable-routing scheme, key anchor-point layout, and sensor configuration are carried out. To compensate for the enveloping-length variation in planar cable routing and simplify the mechanical design, a dynamic cable-exit concept is proposed in the mechanism design stage. Although this concept was not implemented in the second-generation prototype used for force-control validation, and the third-generation design was completed at the system level without being used as the final test platform, this does not affect the validation of the core control algorithms presented in this thesis. In addition, an innovative cable-wrapping method is adopted on the moving platform to achieve continuous multi-turn rotation. A tension-sensor layout is further introduced to provide hardware support for subsequent closed-loop tension control and experimental verification.

For control-oriented modeling, the kinematic relationships among cable lengths, platform pose, and actuator-side states are established. A forward kinematics method based on Jacobian iterative solution is employed to estimate the end-effector pose, and the position-control procedure is verified on the prototype platform. Theoretical analysis and simulation results show that conventional position control based on discretized inverse kinematics can realize basic trajectory tracking; however, under the combined influence of trajectory discretization, cable flexibility, friction, and system coupling, the end-platform acceleration tends to exhibit discontinuity and chattering, which makes it difficult to satisfy the requirement of smooth high-dynamic acceleration simulation. For this reason, the control focus is shifted from position control to force control, with the aim of regulating the generalized force of the platform directly through cable tensions so as to obtain more continuous acceleration output.

On the control side, analysis, simulation, and optimization are carried out for the tension-distribution problem. Since a desired end-platform acceleration in the planar four-cable redundantly actuated mechanism generally corresponds to multiple feasible tension solutions, the effects of positive-tension constraints, load-balancing requirements, and peak-load limitations on the distribution results are investigated. Meanwhile, simulations of end-platform acceleration and theoretical cable tensions are performed along the generated trajectories, which provide the basis for controller design. The study further includes the construction of a force-control experimental platform, single-motor closed-loop torque-control validation, four-motor coordinated tension-control analysis, and end-platform acceleration-tracking experiments.

The theoretical analysis, simulation results, and staged experimental validation indicate that the proposed high-dynamic continuously rotating planar cable-driven parallel robot provides a reasonable foundation for research on continuous rotation and high-dynamic control. Replacing pure position control with force control is shown to be an effective technical route for improving acceleration continuity and dynamic performance. The work presented in this thesis, including planar four-cable mechanism design, kinematic modeling, position-control verification, tension-distribution analysis, force-control platform development, single-motor closed-loop torque control, coordinated multi-cable tension control, and acceleration-tracking experiments, can provide useful references for the design and control of continuously rotating cable-driven parallel robots, and also offer theoretical and engineering support for equipment development in astronaut training, flight simulation, and other high-dynamic motion-simulation scenarios.

\end{eabstract}



